\chapter{切空间、流与 Lie 导数} \label{chap:math-manifolds}

上一部分的谱理论始终在一个固定 Hilbert 空间中比较向量、算子与谱投影；进入流形后，这个默认前提首先失效：不同点上的几何对象属于不同的线性空间，不能未经结构选择就直接相减。在 $\RR^n$ 中，平移提供了这种隐含识别，一般流形却没有预先给定的平移结构。因此本章从光滑结构本身出发，依次建立切空间与微分形式、光滑映射的推前与拉回、向量场生成的流，最后把 Lie 导数定义成“沿流搬回同一点以后的一阶变化率”。联络、度量和 Hodge 结构都在这一基础上逐层加入。

\section{从相容坐标图构造光滑流形}
在平面上，位置可以由一对实数表示。球面上却没有一对处处有效的经纬度：
经度在极点失去意义，但球面本身没有裂口。因此需要区分一个点与给它编号的坐标。
我们先说明“局部区域”和“坐标变换”是什么意思，再把求导运算搬到曲面上。

\section{开集和连续性先于坐标图}
\label{sec:math-manifold-topology}

\begin{definition}[拓扑、邻域与闭集]
集合 \(M\) 上的拓扑 \(\mathcal T\) 是一族子集，含
\(\varnothing,M\)，对任意并和有限交封闭。
\(\mathcal T\) 中的集合称为开集；含点 \(p\) 的开集称为 \(p\) 的开邻域；
补集为开集的集合称为闭集。
\end{definition}

\(\mathbb R^n\) 的通常拓扑取所有这样的集合：集合中每个点都能找到一个
完全包含于集合的小开球。于是开区间是 \(\mathbb R\) 的开集，
\([0,1]\) 是闭集，\([0,1)\) 两者都不是。
这个定义只告诉我们哪些点可在局部一起讨论，尚未赋予距离。

\begin{definition}[连续映射与同胚]
拓扑空间间的映射 \(F:M\to N\) 连续，若 \(N\) 中每个开集 \(V\)
的原像 \(F^{-1}(V)\) 都在 \(M\) 中开。
双射 \(F\) 若它和 \(F^{-1}\) 都连续，称为同胚。
\end{definition}

在欧氏空间，这与高等数学的 \(\varepsilon\)–\(\delta\) 连续性等价。
例如 \(\mathbb R\) 上的 \(x\mapsto x^3\) 是同胚；
\(x\mapsto x^2\) 在全体实数上不是，因为它不单射。
坐标图要求同胚，原因是反向也必须能从坐标唯一、连续地找回原点。

\begin{definition}[两个分离性条件]
若不同点 \(p,q\) 总有互不相交的开邻域，称 \(M\) 为 Hausdorff 空间。
若存在可数个开集 \(B_1,B_2,\ldots\)，使任意开集都是其中若干项的并，
称 \(M\) 第二可数；\(\{B_j\}\) 称为可数基。
\end{definition}

Hausdorff 条件保证极限唯一：若同一序列同时收敛到 \(p\ne q\)，
最终项必须同时落入两个不相交邻域，矛盾。
第二可数也不是空洞的限制。不可数集合若取离散拓扑，
每个单点都开，任何基必须含全部单点，因而不第二可数。
普通 \(\mathbb R^n\) 却有可数基：中心坐标和半径均取有理数的开球。
若 \(M\subset\mathbb R^N\) 已是几何子集，就取子空间拓扑：
\(U\subset M\) 开，当且仅当存在 \(\mathbb R^N\) 的开集 \(O\)
使 \(U=M\cap O\)。这使球面上“删去一个点后的区域”
有确切的开集意义。

\section{两张图怎样描述同一个球面}

\begin{definition}[坐标图与光滑相容]
\(n\) 维坐标图是 \(M\) 的开集 \(U\) 与同胚
\(\varphi:U\to\varphi(U)\subseteq\mathbb R^n\)，其中右边为开集。
两图 \((U,\varphi),(V,\psi)\) 在重叠处光滑相容，若
\[
\psi\circ\varphi^{-1}:
\varphi(U\cap V)\longrightarrow\psi(U\cap V)
\]
及其逆映射都具有任意阶连续偏导数。
覆盖 \(M\) 且两两相容的图族称为光滑图册。
\end{definition}

要求过渡映射光滑比仅要求每张图连续更强。
在实数轴上，\(\varphi(x)=x\) 与
\(\psi(x)=\sqrt[3]{x}\) 都是到实数轴的同胚，
但 \(\psi\circ\varphi^{-1}(u)=\sqrt[3]{u}\)
在零点不可微。这两张图不能同时放进同一光滑图册。
因此“有局部坐标”与“坐标间可无歧义地求导”是两层不同要求。

取单位球面 \(S^2=\{(x,y,z):x^2+y^2+z^2=1\}\)。
删去北极 \(N=(0,0,1)\) 后，沿北极连线投影到赤道平面，得到
\[
\varphi_N(x,y,z)=\left(\frac{x}{1-z},\frac{y}{1-z}\right).
\]
删去南极 \(S=(0,0,-1)\) 后同样得到
\(\varphi_S(x,y,z)=(x/(1+z),y/(1+z))\)。
两张图都覆盖除相应极点外的整个球面。在重叠区域写
\(u=\varphi_N(p)\)，令 \(r^2=|u|^2\)。由球面方程
\(x^2+y^2=1-z^2\) 得
\(r^2=(1+z)/(1-z)\)，所以
\[
\varphi_N^{-1}(u)=
\left(\frac{2u_1}{1+r^2},
\frac{2u_2}{1+r^2},
\frac{r^2-1}{1+r^2}\right).
\]
分母处处非零；这个显式逆映射与 \(\varphi_N\) 都连续，
因而它确实是坐标图，而不仅是一个看起来合理的投影公式。
在两图重叠处 \(u\ne0\)，代入南极图便得到
\[
(\varphi_S\circ\varphi_N^{-1})(u)=\frac{u}{|u|^2}.
\]
它在 \(u\ne0\) 上光滑，且再作用一次便回到 \(u\)。
球面上的点没有因更换坐标而改变；改变的只是它在两张图中的数对。

\begin{definition}[光滑流形与光滑映射]
Hausdorff、第二可数且带有光滑图册的空间称光滑流形。
两个光滑流形间的映射 \(F:M\to N\) 称光滑，
若在任意相容坐标图中，其局部表达
\(\psi\circ F\circ\varphi^{-1}\) 为通常意义下的光滑映射。
\end{definition}

这个定义不要求流形预先嵌入欧氏空间。
球面借助嵌入帮助想象，但时空只需一组相容局部坐标。
对光滑性必须同时查看输入和输出的图；否则一个公式仅说明坐标分量，
尚未说明它定义了哪两个流形之间的映射。

球面上的位置和切向量已经能够不依赖单张坐标图来定义。
可是位于两点的切向量属于不同的空间，单靠坐标图无法定义它们之差；
这正是后续引入联络时需要解决的问题。


\section{切空间、余切空间与外微分}

\begin{workedexample}[圆柱面的切向量]
圆柱面由 \(F(x,y,z)=x^2+y^2-1=0\) 给出。取局部参数化 \(X(\theta,z)=(\cos\theta,\sin\theta,z)\)，求两个坐标切向量并验证它们属于约束的核。

\textbf{解。} 直接求导得
\[
X_\theta=(-\sin\theta,\cos\theta,0),\qquad X_z=(0,0,1).
\]
由于 \(dF_p(v)=2xv_x+2yv_y\)，代入 \(p=X(\theta,z)\) 可得 \(dF_p(X_\theta)=dF_p(X_z)=0\)。两向量线性无关，而正则约束 \(dF_p\ne0\) 的核维数为二，故它们张成整个切空间。\(\theta\) 只能在局部区域取连续实值：绕圆柱一圈会让角度增加 \(2\pi\)，所以单个实值角坐标不能覆盖整个圆柱。
\end{workedexample}

设 $M$ 为光滑 $n$ 维流形，$p\in M$。切向量不应依赖某组坐标来定义。把 $p$ 处切向量定义为作用在光滑函数芽上的导子最直接：线性映射
\[
v:C^\infty(M)\to\RR
\]
若满足 Leibniz 规则
\[
v(fg)=v(f)g(p)+f(p)v(g),
\]
便称为 $p$ 处切向量。所有切向量组成 $n$ 维向量空间 $T_pM$。若 $(x^1,\ldots,x^n)$ 是 $p$ 附近的局部坐标，则
\[
\partial_i\big|_p:=\frac{\partial}{\partial x^i}\Big|_p
\]
构成 $T_pM$ 的一组基，每个 $v\in T_pM$ 唯一写成
\[
v=v^i\partial_i\big|_p.
\]
这里 $v$ 是与坐标无关的几何对象，只有分量 $v^i$ 随坐标选择改变。

切向量还有一个与导子定义等价、在几何直觉上更直接的描述。设 $\gamma:(-\epsilon,\epsilon)\to M$ 是过 $p$ 的光滑曲线，$\gamma(0)=p$。定义其速度对函数的作用
\[
v_\gamma(f)
:=\left.\frac{\dd}{\dd t}\right|_{t=0}f(\gamma(t)).
\]
它满足 Leibniz 规则，因此给出 $T_pM$ 中的导子。若两条曲线 $\gamma_1,\gamma_2$ 对所有 $f\in C^\infty(M)$ 都给出同一导子，就称它们在 $p$ 处一阶等价。于是
\[
T_pM
\cong
\{\text{过 }p\text{ 的光滑曲线}\}/\text{一阶等价}.
\]
在坐标 $x=(x^1,\ldots,x^n)$ 中，
\[
v_\gamma
=\dot x^i(0)\,\partial_i|_p,
\qquad
\dot x^i(0)=\left.\frac{\dd}{\dd t}\right|_0x^i(\gamma(t)).
\]
因此“导子定义”和“曲线速度定义”不是两个不同切空间，而是同一对象的代数与动力学实现。

例：$\RR^n$ 中 $\gamma(t)=p+tv$ 直接给出 $v_\gamma=v^i\partial_i$。由此立得：$T_pM$ 在坐标变换下分量按 Jacobi 矩阵协变变换（$v^a=(\partial y^a/\partial x^i)v^i$），维数恰为 $\dim M=n$；余切空间 $T_p^*M$ 由 $\dd x^i$ 张成并满足 $\dd f(v)=v(f)$；局部平凡化 $TM\cong U\times\RR^n$ 的结构群即 $GL(n,\RR)$。此定义可直接推广到 Banach 流形，场构型 $\phi$ 的变分 $\delta\phi$ 正是构形空间上的切向量。

把各点切空间拼接起来得到切丛
\[
TM=\bigsqcup_{p\in M}T_pM,
\]
它本身是一个 $2n$ 维光滑流形。局部坐标 $(x^i)$ 诱导纤维坐标 $(x^i,v^i)$。余切丛
\[
T^*M=\bigsqcup_{p\in M}T_p^*M
\]
同理具有坐标 $(x^i,p_i)$。后续的向量场、微分形式和张量场本质上都是相应向量丛的光滑截面，因此“场”首先是一个 丛-theoretic 对象，而不是一组坐标分量函数。

余切空间 $T_p^*M$ 是 $T_pM$ 的线性对偶。坐标基 $\{\partial_i\}$ 的对偶基记作 $\{\dd x^i\}$，满足
\[
\dd x^i(\partial_j)=\delta^i{}_j.
\]
对函数 $f\in C^\infty(M)$，它的微分 $\dd f\in T_p^*M$ 由
\begin{equation}\label{eq:math-dg-differential-function}
(\dd f)_p(v):=v(f)
\end{equation}
定义，所以局部表达为
\[
\dd f=\partial_i f\,\dd x^i.
\]
这一定义只使用函数对切向量的方向导数，不需要度量把向量与余向量互相识别。
反对称协变张量组成外代数。记 $\Omega^k(M)=\Gamma(\Lambda^kT^*M)$ 为光滑 $k$-形式空间。若 $\alpha\in\Omega^k(M)$、$\beta\in\Omega^\ell(M)$，外积满足
\[
\alpha\wedge\beta=(-1)^{k\ell}\beta\wedge\alpha.
\]
外微分是唯一满足下列三条性质的线性算子
\begin{equation*}\dd:\Omega^k(M)\longrightarrow\Omega^{k+1}(M):
\qquad
\dd^2=0,
\end{equation*}
并且在函数上与\eqnrefs{eq:math-dg-differential-function}一致，同时满足分次 Leibniz 规则
\[
\dd(\alpha\wedge\beta)
=\dd\alpha\wedge\beta+(-1)^k\alpha\wedge\dd\beta.
\]
在局部坐标中，若
\[
\alpha=\frac{1}{k!}\alpha_{i_1\cdots i_k}
\dd x^{i_1}\wedge\cdots\wedge\dd x^{i_k},
\]
则
\[
\dd\alpha
=\frac{1}{k!}\partial_j\alpha_{i_1\cdots i_k}
\dd x^j\wedge\dd x^{i_1}\wedge\cdots\wedge\dd x^{i_k}.
\]

外微分把所有微分形式组成一个 上链 复
\begin{equation*}0\longrightarrow\Omega^0(M)
\overset{\dd}{\longrightarrow}\Omega^1(M)
\overset{\dd}{\longrightarrow}\cdots
\overset{\dd}{\longrightarrow}\Omega^n(M)
\longrightarrow0,\end{equation*}
因为 $\dd^2=0$。于是可定义 de Rham 上同调
\begin{equation*}H^k_{\mathrm{dR}}(M)
:=\frac{\ker(\dd:\Omega^k\to\Omega^{k+1})}
{\operatorname{im}(\dd:\Omega^{k-1}\to\Omega^k)}.\end{equation*}
闭形式代表局部守恒约束，恰当形式代表可以由低一阶势产生的对象；二者之差由全局拓扑测量。

在 star-shaped 开集 $U\subset\RR^n$ 上，Poincar\'e 引理说明
\[
\dd\alpha=0,
\quad k\ge1
\qquad\Longrightarrow\qquad
\alpha=\dd\beta
\]
局部成立。一个显式 同伦 算符 是
\begin{equation}
(K\alpha)_x(v_1,\ldots,v_{k-1})
=\int_0^1 t^{k-1}
\alpha_{tx}(x,v_1,\ldots,v_{k-1})\,\dd t,
\label{eq:poincare-homotopy-operator}
\end{equation}
这里以 欧氏 坐标识别各点的切空间，收缩向量是固定的 $x$；若用 Euler 场 $\mathcal R_y=y^i\partial_i$，等价公式是 $K\alpha=\int_0^1F_t^*(\iota_{\mathcal R}\alpha)\,\dd t/t$，其中 $F_t(x)=tx$。
并满足 同伦 恒等式
\begin{equation}
\dd K+K\dd=I-\operatorname{ev}_0^*.
\label{eq:poincare-homotopy-identity}
\end{equation}
对正次数形式，$\operatorname{ev}_0^*=0$，所以闭形式自动恰当。

若 $M$ 可定向并带边界，则 Stokes 定理把外微分与边界配对：
\begin{equation}
\int_M\dd\alpha=\int_{\partial M}\alpha,
\qquad
\alpha\in\Omega_c^{n-1}(M).
\label{eq:stokes-general}
\end{equation}
Green、Gauss 和经典 Stokes 公式都是它在不同次数、不同度量识别下的坐标化版本。\eqnrefs{eq:stokes-general} 本身不需要度量；积分由取向与顶次形式定义。

给定向量场 $X\in\Gamma(TM)$，插入算子
\begin{equation*}\iota_{X}:\Omega^k(M)\longrightarrow\Omega^{k-1}(M)
\end{equation*}
定义为
\[
(\iota_{X}\alpha)(X_2,\ldots,X_k)
=\alpha(X,X_2,\ldots,X_k).
\]
它是次数 $-1$ 的反导子：
\[
\iota_{X}(\alpha\wedge\beta)
=(\iota_{X}\alpha)\wedge\beta
+(-1)^k\alpha\wedge(\iota_{X}\beta).
\]
外微分与插入算子都是仅由光滑结构和外代数得到的对象；Lie 导数在本章后半部分会由它们重建。

\section{光滑映射、流与 Lie 括号}
若 $F:M\to N$ 为光滑映射，它把点送到点，因此切向量的自然搬运方向也是从 $M$ 到 $N$。推前
\begin{equation*}F_{*p}:T_pM\to T_{F(p)}N
\end{equation*}
由
\[
(F_{*p}v)(\varphi)=v(\varphi\circ F),
\qquad \varphi\in C^\infty(N),
\]
定义。余向量和微分形式则沿相反方向拉回：对 $\alpha\in\Omega^k(N)$，
\begin{equation*}(F^*\alpha)_p(v_1,\ldots,v_k)
:=\alpha_{F(p)}(F_{*p}v_1,\ldots,F_{*p}v_k).
\end{equation*}
由定义直接可验证
\begin{equation*}F^*(\dd\alpha)=\dd(F^*\alpha),
\qquad
F^*(\alpha\wedge\beta)=F^*\alpha\wedge F^*\beta.
\end{equation*}
这两条自然性关系是以后几乎所有外形式计算能够保持坐标无关的原因。

光滑映射的微分不仅给出向量的搬运，还控制局部几何形状。若 $F:M^m\to N^n$ 在 $p$ 附近秩恒为 $r$，constant 秩 定理 保证存在局部坐标，使
\[
F(x^1,\ldots,x^m)
=(x^1,\ldots,x^r,0,\ldots,0).
\]
因此三个基本结论统一为同一正规形：
\begin{itemize}
\item 若 $m=n=r$，$F$ 局部是微分同胚，这就是 逆 函数 定理；
\item 若 $r=m\le n$，$F$ 是局部 浸入；
\item 若 $r=n\le m$，$F$ 是局部 淹没。
\end{itemize}
若 $y\in N$ 是 正则 值，即对每个 $p\in F^{-1}(y)$ 都有 $F_{*p}$ 满射，则
\begin{equation*}F^{-1}(y)\subset M\end{equation*}
是维数 $m-n$ 的嵌入子流形，而且
\begin{equation}
T_pF^{-1}(y)=\ker F_{*p}.
\label{eq:tangent-regular-level-set}
\end{equation}
这给出了“约束方程的切空间”最一般的几何表达。

三个典型 水平集 立即说明此公式。球面 $S^{n-1}=F^{-1}(1)$（$F(x)=|x|^2$，$\dd F_x(v)=2x\cdot v$ 在 $|x|=1$ 时满射），$T_xS^{n-1}=\{v\in\RR^n:x\cdot v=0\}$ 是与位置向量正交的超平面；正交群 $O(n)=F^{-1}(I)$（$F(A)=A^TA$，$\dd F_A(B)=A^TB+B^TA$ 满射），$T_IO(n)=\mathfrak{so}(n)$ 即反对称矩阵 Lie 代数；特殊线性群 $SL(n,\RR)=F^{-1}(1)$（$F(A)=\det A$），$T_ISL(n,\RR)=\mathfrak{sl}(n,\RR)$ 即迹零矩阵 Lie 代数。正则 水平集 定理是约束系统力学的基础：约束 $F(q)=0$（$F:Q\to\RR^k$）在 正则 条件下给出约束子流形 $Q_{\mathrm c}=F^{-1}(0)$，切空间是虚位移空间，Lagrange 乘子正是约束力；微正则系综（$H=E$）与规范约束（$\nabla\cdot E=\rho$）中的'约束即光滑子流形'是同一图景。

例：$F:S^2\to\RR^3$ 为包含映射、$\alpha=x\,\dd y-y\,\dd x$ 时，$\dd\alpha=2\,\dd x\wedge\dd y$，$F^*(\dd\alpha)$ 恰为球面面积形式的 2 倍。拉回与外微分交换是 de Rham 上同调自然性的基础：$F^*$ 诱导同态 $F^*:H^k_{\mathrm{dR}}(N)\to H^k_{\mathrm{dR}}(M)$。它保证坐标不变的物理结论不随拉回改变，例如电磁学中场强 $F=\dd A$ 在规范变换 $A\mapsto A+\dd\chi$ 下不变，正是 $\dd^2=0$ 与拉回交换性的结合。

正则 值 定理 只讨论目标子流形退化为一个点的情形。更一般地，设 $F:M\to N$ 光滑，$S\subset N$ 为嵌入子流形。若对每个 $p\in F^{-1}(S)$ 都有
\begin{equation}
{
\dd F_p(T_pM)+T_{F(p)}S=T_{F(p)}N,
}
\label{eq:math-dg-transversality}
\end{equation}
则称 $F$ 与 $S$ 横截，记为
\[
F\pitchfork S.
\]
这个条件允许 $\dd F_p$ 本身不是满射；只要求它缺失的方向由 $S$ 的切空间补齐。

\begin{theorem}[原像定理]\label{thm:math-dg-preimage-transverse}
若 $F\pitchfork S$，则 $F^{-1}(S)$ 是 $M$ 的嵌入子流形，并且
\begin{equation}
\operatorname{codim}_M F^{-1}(S)
=\operatorname{codim}_N S.
\label{eq:math-dg-transverse-codimension}
\end{equation}
其切空间为
\begin{equation}
T_pF^{-1}(S)
=\{v\in T_pM:\dd F_p(v)\in T_{F(p)}S\}.
\label{eq:math-dg-transverse-tangent-space}
\end{equation}
\end{theorem}

例：$F(t)=(t,t^2)$ 与 $x$ 轴在 $t=0$ 处横截，$F^{-1}(S)=\{0\}$ 为零维子流形；$F(t)=(t,0)$ 非横截，$F^{-1}(S)=\RR$ 为一维，微扰 $F_\varepsilon(t)=(t,\varepsilon)$ 后横截 原像 为空，体现横截性的稳定性。约束力学中 $Q_{\mathrm c}=\{q:F(q)=0\}$（$F:\RR^{3N}\to\RR^k$）在 $\dd F_q$ 满秩时为 $3N-k$ 维子流形；纤维积 $M\times_NP$ 即 $(F,G)$ 与对角 $\Delta_N$ 的横截 原像，维数 $\dim M+\dim P-\dim N$。一般地，两 $a,b$ 维子流形横截相交给出 $\dim(A\cap B)=a+b-n$。横截性也是模空间理论的基础：瞬子模空间 $\mathcal M$ 由自对偶方程 $F_A^+=0$ 定义，横截性保证光滑，维数由 Atiyah--Singer 指标定理给出；其一般形式正是 Sard 定理——横截相交在参数扰动下保持稳定。

横截性之所以特别重要，是因为它是稳定且泛型的几何条件。Sard 定理说明，对光滑映射 $F:M^m\to N^n$，临界值 的集合在 $N$ 中具有测度零。于是“几乎每个” $y\in N$ 都是 正则 值；对这样的 $y$，水平集 $F^{-1}(y)$ 自动是光滑子流形。更强的 横截性 定理 说明，在合适函数空间拓扑下，横截映射形成稠密集合：任意映射通常只需任意小扰动，就能把非横截交会变成横截交会。

这给出“维数计数”的严格来源。设 $A,B\subset N^n$ 是维数分别为 $a,b$ 的嵌入子流形。若它们在交点横截，
\[
T_pA+T_pB=T_pN,
\]
则
\begin{equation*}\dim(A\cap B)=a+b-n.\end{equation*}
若右边为负，泛型情形下交集为空；若为零，泛型交集是离散点。这就是 differential 拓扑 中“未知量维数减独立约束数”的精确版本。

更一般地，给定 $F:M^m\to N^n$ 与 $G:P^p\to N^n$，若映射
\[
F\times G:M\times P\to N\times N
\]
与对角子流形
\[
\Delta_N:=\{(y,y):y\in N\}
\]
横截，则 纤维积
\begin{equation*}M\times_NP
:=\{(x,z):F(x)=G(z)\}\end{equation*}
是维数
\[
m+p-n
\]
的子流形。这个构造以后在 规范 匹配、边界条件拼接、模 空间 和统计约束模型里都会反复出现。

\begin{exbox}[独立约束与 Lagrange 乘子 的几何前提]
对独立约束 $F=(F^1,\ldots,F^k)=0$，在 $\dd F$ 满秩的假设下求约束流形的切空间，并推出受限极值的 Lagrange 乘子条件。

\textbf{解。} 设
\[
F=(F^1,\ldots,F^k):\RR^n\to\RR^k,
\]
约束集为
\[
\mathcal C=F^{-1}(0).
\]
若 $\dd F_x$ 在所有 $x\in\mathcal C$ 上满秩 $k$，则 $0$ 是 正则 值，故 $\mathcal C$ 是 $n-k$ 维子流形，并且
\[
T_x\mathcal C=\ker\dd F_x.
\]
临界点条件
\[
\dd f_x|_{T_x\mathcal C}=0
\]
等价于 $\dd f_x$ 落在 $\ker\dd F_x$ 的 annihilator 中，也就是由 $\dd F^1_x,\ldots,\dd F^k_x$ 张成。因此存在 Lagrange 乘子 $\lambda_a$ 使
\[
\dd f_x=\sum_{a=1}^k\lambda_a\dd F^a_x.
\]
所以 Lagrange 乘子 法成立的核心并不是代数消元，而是约束映射在解集上具有满秩，从而约束集真的是光滑子流形。
\end{exbox}
向量场把“静态的切向量”提升为点的运动。设 $X\in\Gamma(TM)$，局部流 $\Phi_t$ 由初值问题
\begin{equation}\label{eq:math-dg-flow-ode}
\frac{\dd}{\dd t}\Phi_t(p)=X_{\Phi_t(p)},
\qquad
\Phi_0(p)=p
\end{equation}
定义。在共同存在的时间区间内，唯一性给出
\[
\Phi_{t+s}=\Phi_t\circ\Phi_s,
\qquad
\Phi_t^{-1}=\Phi_{-t}.
\]
因此流不仅描述点如何运动，还提供了把几何对象从 $\Phi_t(p)$ 搬回 $p$ 的微分同胚。这正是 Lie 导数所需的比较机制。

局部流的存在唯一性直接继承自常微分方程理论：在任一坐标图中，\eqnrefs{eq:math-dg-flow-ode} 是
\[
\dot x^i(t)=X^i(x(t)).
\]
只要 $X$ 光滑，就在每个初值附近存在唯一 极大 积分曲线。若所有积分曲线都能延伸到全部 $t\in\RR$，则称 $X$ 完全，此时 $\Phi_t$ 对所有时间都定义，并形成真正的一参数微分同胚群
\[
\Phi:\RR\times M\to M,
\qquad
\Phi_{t+s}=\Phi_t\circ\Phi_s.
\]
紧流形上的任意光滑向量场都 完全；非紧流形则可能有限时间逃向无穷远。例如 $\RR$ 上
\[
X=x^2\partial_x
\]
的积分曲线满足
\[
x(t)=\frac{x_0}{1-tx_0},
\]
对 $x_0>0$ 在 $t=1/x_0$ 爆破，因此不是 完全。

若 $F:M\to N$ 是微分同胚，向量场自然推前为
\[
(F_*X)_{F(p)}=F_{*p}X_p.
\]
若 $X$ 与 $Y$ 分别是 $M,N$ 上向量场，并满足
\[
F_*X=Y\circ F,
\]
则称二者 $F$-related。其流也满足
\begin{equation}
F\circ\Phi_t^X=\Phi_t^Y\circ F
\label{eq:related-flow-naturality}
\end{equation}
在共同存在区间内成立。这种自然性将在 Lie 导数、联络拉回与对称性分析中反复使用。

例：$F(x)=Ax$（$A$ 可逆）把径向场 $X=x^i\partial_i$ 推前为同方向径向场，$\ee^{tI}$ 与 $A$ 对易即给出流共轭；极坐标 $F:(r,\theta)\mapsto(r\cos\theta,r\sin\theta)$ 把 $\partial_\theta$ 推前为旋转场 $-y\partial_x+x\partial_y$。此共轭性是对称性约化的基础：$G$ 作用下不变向量场推前到商空间，流的共轭保证约化流与原流相容；Noether 定理即对称流与 Hamilton 流相容 $[X,X_H]=0$ 时给出守恒量。

两个向量场 $X,Y$ 的 Lie 括号定义为导子的交换子
\begin{equation*}[X,Y](f):=X(Yf)-Y(Xf).
\end{equation*}
由于右端仍满足 Leibniz 规则，$[X,Y]$ 确实是向量场。若
\[
X=X^i\partial_i,
\qquad
Y=Y^i\partial_i,
\]
则
\[
[X,Y]
=\bigl(X^j\partial_jY^i-Y^j\partial_jX^i\bigr)\partial_i.
\]
这个式子只含一阶导数，但它测量的是两个无穷小流的二阶不交换效应。

例：$X=\partial_x$、$Y=x\partial_y$ 时 $[X,Y]=\partial_y$，群交换子 $\Psi_{-s}\Phi_{-t}\Psi_s\Phi_t(x,y)=(x,y+st)$ 的位移 $st$ 恰为括号的流；Heisenberg 代数 $X=\partial_x-\frac12y\partial_z$、$Y=\partial_y+\frac12x\partial_z$ 给出 $[X,Y]=\partial_z$，量子上对应正则对易 $[\hat x,\hat p]=\ii\hbar$；$SO(3)$ 的两个旋转场括号即交叉积法则，是刚体 Euler 方程的几何来源。左不变向量场在恒等元处的括号给出 $\mathfrak g$，群交换子在恒等元附近展开为 $ts[X,Y]$；Hamilton 向量场满足 $[X_H,X_F]=X_{\{H,F\}}$，同一括号在量子层面对应算子交换子、在规范理论中对应场强的非交换部分 $[A_\mu,A_\nu]$。

\section{Lie 导数与 Cartan 恒等式}
流已经给出了跨底点搬运对象的方法。对任意张量场 $T$，Lie 导数定义为沿 $X$ 的流先拉回再求变化率：
\begin{equation}\label{eq:math-dg-synthesis-lie-flow}
\mathcal L_XT
:=\left.\frac{\dd}{\dd t}\right|_{t=0}\Phi_t^*T.
\end{equation}
这里对反变指标的“拉回”使用 $\Phi_{-t*}$ 把向量从 $\Phi_t(p)$ 搬回 $p$。定义的要点是比较永远发生在同一底点，因此不需要先选联络。

对函数与向量场，\eqnrefs{eq:math-dg-synthesis-lie-flow}分别给出
\begin{equation*}\mathcal L_Xf=X(f),
\qquad
\mathcal L_XY=[X,Y].
\end{equation*}
对一形式 $\alpha$，利用配对 $\alpha(Y)$ 的 Leibniz 规则得到
\[
(\mathcal L_X\alpha)(Y)
=X\bigl(\alpha(Y)\bigr)-\alpha([X,Y]).
\]
对一般 $(r,s)$ 张量，Lie 导数与张量积和缩并相容；在坐标基中可写成
\begin{align*}
(\mathcal L_XT)^{i_1\cdots i_r}{}_{j_1\cdots j_s}
={}&X^k\partial_kT^{i_1\cdots i_r}{}_{j_1\cdots j_s}
-\sum_{a=1}^r
T^{i_1\cdots k\cdots i_r}{}_{j_1\cdots j_s}
\partial_kX^{i_a}\\
&+\sum_{b=1}^s
T^{i_1\cdots i_r}{}_{j_1\cdots k\cdots j_s}
\partial_{j_b}X^k.
\end{align*}
这只是坐标表达，不是定义本身。
外微分与任意拉回交换，因此对流拉回求导后立即得到
\begin{equation}\label{eq:math-dg-synthesis-lie-d-commute}
[\mathcal L_X,\dd]=0.
\end{equation}
对微分形式还存在更强的无坐标表达。

对次数齐次算子 $A,B$，定义 分次 对易子
\[
[A,B]_{\rm gr}
:=AB-(-1)^{|A||B|}BA.
\]
在微分形式代数上，$\dd$ 的次数为 $+1$，$\iota_X$ 的次数为 $-1$，$\mathcal L_X$ 的次数为 $0$。Cartan 恒等式可以压缩为
\[
\mathcal L_X=[\dd,\iota_X]_{\rm gr}.
\]
进一步有整套 Cartan 演算：
\begin{align}
[\mathcal L_X,\mathcal L_Y]&=\mathcal L_{[X,Y]},
\label{eq:cartan-lie-lie}\\
[\mathcal L_X,\iota_Y]&=\iota_{[X,Y]},
\label{eq:cartan-lie-iota}\\
[\iota_X,\iota_Y]_{\rm gr}&=0,
\label{eq:cartan-iota-iota}\\
[\mathcal L_X,\dd]&=0.
\label{eq:cartan-lie-d}
\end{align}
这些关系说明 $\Omega^\bullet(M)$ 不只是一个微分分次代数，而且携带由向量场 Lie 代数 作用产生的完整微分代数结构。

Cartan 恒等式把“沿流的变化”完全改写为外代数操作。它的优势在于只依赖 $X$、$\dd$ 与插入算子，不需要坐标、度量或联络。后面出现体积形式、联络形式以及几何输运方程时，都将继续使用这一同一公式，而不再重新定义 Lie 导数。

若 $M$ 可定向并选定一个处处非零的体积形式 $\mu\in\Omega^n(M)$，则向量场 $X$ 关于 $\mu$ 的散度由
\begin{equation*}\mathcal L_X\mu=(\operatorname{div}_\mu X)\mu\end{equation*}
定义。因为最高次形式自动满足 $\dd\mu=0$，Cartan 公式给出
\begin{equation*}(\operatorname{div}_\mu X)\mu
=\dd(\iota_X\mu).\end{equation*}
于是对紧支撑 $X$，Stokes 定理立即推出
\[
\int_M\operatorname{div}_\mu X\,\mu=0
\]
在无边界情形成立；有边界时则产生通量项。这是 散度 定理 的坐标无关形式。

若 $\rho_t$ 是随流输运的密度，使总质量形式 $\rho_t\mu$ 满足
\[
\partial_t(\rho_t\mu)+\mathcal L_X(\rho_t\mu)=0,
\]
则由 Leibniz 规则得到连续性方程
\begin{equation*}\partial_t\rho+X(\rho)+\rho\operatorname{div}_\mu X=0.\end{equation*}
在 欧氏 体积形式下就是
\[
\partial_t\rho+\nabla\cdot(\rho X)=0.
\]
因此守恒律中的“散度结构”首先来自 陀螺-形式 的 Lie 输运，而不依赖笛卡尔坐标。

更一般地，设 $D_t=\Phi_t(D_0)$ 是由 $X$ 的流推动的移动区域，$\alpha_t$ 是时变 $k$-形式。则 Reynolds 输运 定理 的微分形式版本为
\begin{equation*}\frac{\dd}{\dd t}\int_{D_t}\alpha_t
=\int_{D_t}(\partial_t\alpha_t+\mathcal L_X\alpha_t).\end{equation*}
证明只需把积分域拉回固定的 $D_0$：
\[
\int_{D_t}\alpha_t
=\int_{D_0}\Phi_t^*\alpha_t,
\]
再对 $t$ 求导。该公式把流体输运、守恒律和 Lie 导数直接连接起来。

\begin{exbox}[旋转流保持径向一形式]
在平面上取转动向量场 $X=-y\partial_x+x\partial_y$ 与一形式 $\alpha=x\dd x+y\dd y$，求 $\mathcal L_X\alpha$ 并判断转动流是否保持它。

\textbf{解。} 在 $\RR^2$ 上取
\[
X=-y\partial_x+x\partial_y,
\qquad
\alpha=x\,\dd x+y\,\dd y.
\]
因为
\[
\dd\alpha=0,
\qquad
\iota_{X}\alpha=-xy+xy=0,
\]
由\eqnrefs{eq:math-dg-synthesis-cartan}立刻得到
\[
\mathcal L_X\alpha=0.
\]
这里 $\alpha=\tfrac12\dd(x^2+y^2)$，所以结果表达的正是旋转流保持径向距离函数。
\end{exbox}


切向量、微分形式与流都已具有坐标无关的定义。下面利用参数曲线、球面坐标和闭形式的反例，检验推前、外微分与 Lie 括号的实际含义。

\section{曲线速度与点处求导为何是同一对象}
\label{sec:math-tangent-derivations}

球面上的一条曲线 \(\gamma:(-\epsilon,\epsilon)\to S^2\) 有速度，
但把 \(\dot\gamma(0)\) 看作三维向量只适用于已经嵌入
\(\RR^3\) 的曲面。一个不依赖嵌入的办法是记录曲线经过点
\(p=\gamma(0)\) 时，怎样改变所有光滑函数的值。

\begin{definition}[曲线切向量]
经过 \(p\in M\) 的两条可微曲线 \(\gamma_1,\gamma_2\) 等价，
若在一张经过 \(p\) 的坐标图 \(\varphi\) 中有
\((\varphi\circ\gamma_1)'(0)=(\varphi\circ\gamma_2)'(0)\)。
等价类称为 \(p\) 处的切向量，其集合记为 \(T_pM\)。
\end{definition}

这里的“一张”不会导致坐标依赖。若换到图 \(\psi\)，链式法则给
\[
(\psi\circ\gamma)'(0)
=D(\psi\circ\varphi^{-1})_{\varphi(p)}
(\varphi\circ\gamma)'(0).
\]
两曲线在 \(\varphi\) 中速度相同，则在 \(\psi\) 中也相同。
这个等式同时告诉我们：切向量本身不变，分量被坐标变换的 Jacobian 相乘。

\begin{definition}[点处导子]
点 \(p\) 处的导子是一个把 \(p\) 附近光滑实值函数送到实数的线性映射
\(v\)，满足 Leibniz 法则
\[
v(fg)=f(p)v(g)+g(p)v(f).
\]
只在 \(p\) 附近相同的函数视为同一个输入；这种局部函数称为 \(p\) 处的芽。
\end{definition}

曲线 \(\gamma\) 给出导子
\(v_\gamma(f)=\left.\dd f(\gamma(t))/\dd t\right|_{t=0}\)。
乘积求导法则说明它确实满足上面的定义。反过来，导子并不会包含比曲线速度
更多的信息。

\begin{theorem}[曲线定义与导子定义的对应]
在 \(n\) 维光滑流形 \(M\) 上，每个 \(p\) 处的导子都由唯一的曲线切向量给出。
给定坐标 \(x=(x^1,\ldots,x^n)\)，该导子可唯一写成
\[
v=\sum_{i=1}^n v(x^i)\,\partial_i\big|_p,\qquad
\partial_i\big|_p(f)=\partial_i(f\circ x^{-1})(x(p)).
\]
\end{theorem}
\begin{proof}
先由 Leibniz 法则作用于常数函数 \(1=1\cdot1\) 得 \(v(1)=0\)，
故所有常数的导数为零。把局部坐标中的 \(f\) 限制在 \(x(p)\) 附近的小凸球，
沿连接 \(x(p)\) 与 \(x\) 的直线使用微积分基本定理：
\[
f(x)-f(x(p))
=\sum_i(x^i-x^i(p))f_i(x),\qquad
f_i(x)=\int_0^1\partial_i f(x(p)+t(x-x(p)))\,\dd t.
\]
这里 \(f_i\) 光滑且 \(f_i(x(p))=\partial_i f(x(p))\)。
对等式作用 \(v\)，常数项为零，Leibniz 法则使
\((x^i-x^i(p))v(f_i)\) 在 \(p\) 处消失，得到
\[
v(f)=\sum_i v(x^i)\partial_i f(x(p)).
\]
因此数 \(v^i=v(x^i)\) 决定全部作用。取
\(\gamma(t)=x^{-1}(x(p)+t(v^1,\ldots,v^n))\)，小 \(t\) 时该曲线有定义，
链式法则给 \(v_\gamma(f)=v(f)\)。唯一性由各坐标函数的值确定。
\end{proof}

把球面当作 \(\RR^3\) 的子集时，约束
\(x^2+y^2+z^2=1\) 沿曲线求导给
\(p\cdot\dot\gamma(0)=0\)。
反过来，每个满足 \(p\cdot v=0\) 的三维向量都可作为球面上一条曲线
的初速度，例如把 \(p+tv\) 归一化。
因此 \(T_pS^2=\{v\in\RR^3:p\cdot v=0\}\)。
这算出了切空间，也验证抽象定义与熟悉的切平面一致。

\begin{proposition}[正则约束的切空间]
设 \(F:\RR^N\to\RR^m\) 光滑，
\(M=F^{-1}(0)\)，并且 \(DF_p\) 在每个 \(p\in M\)
的秩为 \(m\)。则 \(M\) 在每个点附近是 \((N-m)\) 维
光滑流形，且
\[
T_pM=\ker DF_p.
\]
\end{proposition}
\begin{proof}
由秩条件，可重新排列环境坐标，使 \(DF_p\) 对最后 \(m\)
个坐标的偏导方阵可逆。多变量微积分的隐函数定理遂把这些坐标
局部写成前 \(N-m\) 个坐标的光滑函数，这给出一张局部图。
若 \(\gamma(t)\subset M\) 且 \(\gamma(0)=p\)，
对 \(F(\gamma(t))=0\) 求导，得
\(DF_p\dot\gamma(0)=0\)，故 \(T_pM\subseteq\ker DF_p\)。
局部图的 \(N-m\) 个坐标速度线性独立并全在该核内；
由秩—零度定理，核维数也是 \(N-m\)，故两者相等。
\end{proof}

球面取 \(F(x,y,z)=x^2+y^2+z^2-1\)，
\(DF_pv=2p\cdot v\)，于是回到上面的切平面。
秩条件不能省略：锥面
\(x^2+y^2-z^2=0\) 在顶点的 \(DF_0=0\)，
形式上的核是整个 \(\RR^3\)，却不存在一个二维光滑切平面。
这就是约束方程虽可写下，几何对象却可能在特殊点变坏的具体例子。


\par
\begin{workedexample}
在圆柱面 \(x^2+y^2=1\) 的局部坐标 \((\theta,z)\) 中，
求 \(\partial_\theta,\partial_z\) 作为 \(\RR^3\) 向量的表达，
并验证它们满足正则约束切空间的判据。
解释为何 \(\theta\) 仍不能作为覆盖整个圆柱面的单张实值坐标。
\textbf{解。} 参数化 \(F(\theta,z)=(\cos\theta,\sin\theta,z)\) 给
\[
\partial_\theta F=(-\sin\theta,\cos\theta,0),\qquad
\partial_zF=(0,0,1).
\]
两向量均与约束函数 \(G=x^2+y^2-1\) 的梯度 \(\nabla G=(2x,2y,0)\) 正交，且线性无关，故张成正则约束的二维切空间。角变量以 \(2\pi\) 为周期，不能在整圆周上成为单张实值坐标；需用多个坐标片。
\end{workedexample}
\par
\section{映射怎样推动速度}

若 \(F:M\to N\) 光滑，曲线 \(\gamma\) 经 \(F\) 变成
\(F\circ\gamma\)。由此定义点 \(p\) 处的微分
\[
dF_p:T_pM\to T_{F(p)}N,\qquad
dF_p(v)(h)=v(h\circ F).
\]
它是线性的，且因函数复合满足
\(d(G\circ F)_p=dG_{F(p)}\circ dF_p\)。
在坐标中，\(dF_p\) 正是熟悉的 Jacobian 矩阵；
写清定义域与值域能防止把两个不同点的速度直接相加。

\section{对偶向量与张量}

线性代数中，行向量作用于列向量得到一个数。流形上也一样：
\begin{definition}[余切向量与张量]
\(T_pM\) 上的线性泛函组成余切空间 \(T_p^*M\)。
函数 \(f\) 的微分 \(df_p\) 由 \(df_p(v)=v(f)\) 定义。
一个 \((r,s)\) 型张量是在同一点 \(p\) 处，对 \(r\) 个余切向量和
\(s\) 个切向量分别线性的多重映射，值为实数。
\end{definition}

坐标函数的微分 \(\dd x^i\) 与 \(\partial_j\) 对偶，因为
\(\dd x^i(\partial_j)=\delta^i{}_j\)。
从现在起采用爱因斯坦求和约定：同一乘积项内一个上指标与一个下
指标各出现一次时，自动对该指标的全部取值 \(1,\dots,n\) 求和，
不再写求和号；上例 \(v=\sum_i v(x^i)\partial_i\) 今后就写成
\(v=v(x^i)\partial_i\)。
若新坐标为 \(y^a=y^a(x)\)，上面的链式法则给
\[
\frac{\partial}{\partial y^a}
=\frac{\partial x^i}{\partial y^a}\frac{\partial}{\partial x^i},
\qquad
\dd y^a=\frac{\partial y^a}{\partial x^i}\dd x^i .
\]
两种 Jacobian 互逆，故对偶配对不随坐标改变。
在平面极坐标区域 \(r>0\)，直接微分
\(x=r\cos\theta,y=r\sin\theta\)，得
\[
\partial_r=\cos\theta\,\partial_x+\sin\theta\,\partial_y,
\qquad
\partial_\theta=-r\sin\theta\,\partial_x
+r\cos\theta\,\partial_y.
\]
反过来，由 \(r=\sqrt{x^2+y^2}\) 与局部角函数得到
\[
\dd r=\frac{x\,\dd x+y\,\dd y}{r},
\qquad
\dd \theta=\frac{-y\,\dd x+x\,\dd y}{r^2}.
\]
于是 \(\dd r(\partial_\theta)=0\)、
\(\dd \theta(\partial_\theta)=1\)，确实是一对对偶基。
但角度 \(\theta\) 不能在整个穿孔平面选成连续单值函数；
右侧的公式却拼成一个处处光滑的全局一形式。
下一章将证明它虽处处局部等于角度的微分，却不是全局函数的微分。
例如点处的双线性型 \(b_p(v,w)\) 是 \((0,2)\) 型张量；记
\((\alpha\otimes\beta)(v,w)=\alpha(v)\beta(w)\)，则它可写成
\(b=b_{ij}\dd x^i\otimes \dd x^j\)，这时才出现指标。
因此“张量”首先是多重线性映射，带指标的数组只是选基后的坐标记录。


\section{反对称张量、微分形式与积分}
\label{sec:math-forms-from-tensors}

计算沿曲线的积分时，线元 \(\dd x^i\) 会与函数系数一起变换，
使积分与坐标选择无关。面积积分需要同样的机制，但还要求交换两个方向时
方向反号。这把我们带到反对称张量。

\begin{definition}[微分形式与楔积]
\(p\) 处的 \(k\)-形式是 \(T_pM\) 上的交替 \(k\) 重线性映射：
交换任意两个输入会使值变号。随 \(p\) 光滑变化的 \(k\)-形式记作
\(\Omega^k(M)\)。对一形式 \(\alpha,\beta\)，
\[
(\alpha\wedge\beta)(v,w)
=\alpha(v)\beta(w)-\alpha(w)\beta(v).
\]
若 \(\alpha\) 为 \(k\)-形式、\(\beta\) 为 \(\ell\)-形式，
一般楔积由下面的交替化公式固定归一化：
\[
(\alpha\wedge\beta)(v_1,\ldots,v_{k+\ell})
=\frac1{k!\ell!}\sum_{\pi\in S_{k+\ell}}
\operatorname{sgn}(\pi)\,
\alpha(v_{\pi(1)},\ldots,v_{\pi(k)})
\beta(v_{\pi(k+1)},\ldots,v_{\pi(k+\ell)}).
\]
这里 \(S_{k+\ell}\) 是全部输入次序的排列，
\(\operatorname{sgn}(\pi)\) 在偶数次交换时为正、奇数次时为负。
交换两组输入需要 \(k\ell\) 次换位，因而
\(\alpha\wedge\beta=(-1)^{k\ell}\beta\wedge\alpha\)。
\end{definition}

在坐标中，\(\dd x^1\wedge \dd x^2(\partial_1,\partial_2)=1\)；
交换两个方向则为 \(-1\)。若 \(y=y(x)\)，将
\(\dd y^a=(\partial_i y^a)\dd x^i\) 代入楔积，得到二维情形
\[
\dd y^1\wedge \dd y^2
=\det\!\left(\frac{\partial(y^1,y^2)}
{\partial(x^1,x^2)}\right)\dd x^1\wedge \dd x^2.
\]
雅可比行列式由反对称性自动产生，这就是带方向面积积分的坐标变换律。

函数的微分已由 \(\dd f(v)=v(f)\) 定义。要对一形式继续求导，
就要求新的操作既与坐标变换相容，又保留
\(\dd(\dd f)=0\) 这一混合偏导交换的事实。

\begin{definition}[外微分]
在每张坐标图中，对
\(\alpha=\sum_I a_I\,\dd x^{i_1}\wedge\cdots\wedge \dd x^{i_k}\)，
定义
\[
\dd\alpha=\sum_I \dd a_I\wedge \dd x^{i_1}\wedge\cdots\wedge \dd x^{i_k}.
\]
下面验证这些局部公式拼成同一个
\(\dd:\Omega^k(M)\to\Omega^{k+1}(M)\)。
\end{definition}

上式为何与坐标无关？先在函数上，链式法则保证 \(\dd f\) 的定义不依赖图。
对新坐标 \(y^a\)，写
\(\dd y^a=(\partial_i y^a)\dd x^i\) 并按局部公式求导，得到
\(\dd(\dd y^a)=(\partial_j\partial_i y^a)\dd x^j\wedge \dd x^i=0\)：
二阶偏导对称，楔积反对称。局部公式还直接满足分级 Leibniz 法则
\[
\dd(\alpha\wedge\beta)
=\dd\alpha\wedge\beta+(-1)^k\alpha\wedge \dd\beta
\quad(\alpha\in\Omega^k(M)).
\]
于是把任意形式改写成 \(\dd y^a\) 的楔积后再求导，
与原坐标下的结果相同。相同的对称—反对称抵消也给出
\(\dd^2\alpha=0\)：系数的二阶偏导关于指标对称，
两个 \(\dd x\) 的楔积反对称，两者相乘后逐项相消。

例如平面上的一形式 \(\alpha=x\,\dd y-y\,\dd x\) 满足
\[
\dd\alpha=\dd x\wedge \dd y-\dd y\wedge \dd x=2\,\dd x\wedge \dd y.
\]
在单位圆盘上积分右边得 \(2\pi\)，沿边界逆时针计算左边也得
\(2\pi\)。这既是 Stokes 定理的一个直接检查，也说明边界方向
必须与面积方向配套。此处用的是高数里的 Green 公式；高维的一般形式
\(\int_M\dd\omega=\int_{\partial M}\omega\) 与此一致，本书用到时直接引用。


\par
\begin{workedexample}
在 \(\RR^2\) 取 \(X=x\partial_y\)、\(Y=\partial_x\)。
从定义算 \([X,Y]\)，并说明为何两次沿流移动的顺序有影响。
\textbf{解.} 对任意光滑 \(f\)，
\[
[X,Y]f=x\partial_y\partial_x f-\partial_x(x\partial_yf)=-\partial_yf,
\]
故 \([X,Y]=-\partial_y\)。沿 \(Y\) 先平移再沿 \(X\) 移动，与反过来次序的二阶差由该括号控制；这里 \(X\) 的强度依赖 \(x\)，而 \(Y\) 会改变 \(x\)，所以两流程不交换。
\end{workedexample}
\par
\section{零旋度何时真的意味着存在势函数}

写一形式 \(\alpha=P(x,y)\,\dd x+Q(x,y)\,\dd y\)。
直接用外微分定义，
\(\dd\alpha=(\partial_xQ-\partial_yP)\dd x\wedge \dd y\)。
若 \(\alpha=\dd f\)，则 \(\dd\alpha=\dd^2f=0\)。
反方向却需要检查定义域的形状，而不是只检查偏导。

\begin{definition}[闭形式、恰当形式与星形区域]
满足 \(\dd\alpha=0\) 的形式称为闭形式；
若存在低一阶形式 \(\beta\) 使 \(\alpha=\dd\beta\)，
称 \(\alpha\) 为恰当形式。
平面区域 \(U\) 相对点 \(x_0\) 是星形的，若每个
\(x\in U\) 与 \(x_0\) 的整个线段都在 \(U\) 内。
\end{definition}

\begin{theorem}[一形式的局部势函数（Poincaré 引理）]
若 \(U\subset\RR^2\) 是星形开集，
\(P,Q\) 连续可微，且 \(\partial_xQ=\partial_yP\)，
则存在函数 \(f\) 使 \(\alpha=\dd f\)。
\end{theorem}
\begin{proof}
平移原点到星形中心，写 \(\alpha=\alpha_i(x)\dd x^i\) 并定义
\[
f(x)=\int_0^1 x^i\alpha_i(tx)\,\dd t.
\]
线段始终在 \(U\) 内，故积分有定义。
对第 \(j\) 个坐标求导，闭条件
\(\partial_j\alpha_i=\partial_i\alpha_j\) 给
\[
\partial_j f
=\int_0^1\left[\alpha_j(tx)
+t x^i\partial_j\alpha_i(tx)\right]\dd t
=\int_0^1\frac{\dd}{\dd t}\bigl[t\alpha_j(tx)\bigr]\dd t
=\alpha_j(x).
\]
所以 \(\dd f=\alpha\)。证明同时给出可实际计算的势函数，
而不只断言它存在。
\end{proof}

去掉星形条件就有反例。在穿孔平面
\(U=\RR^2\setminus\{0\}\) 上取
\[
\alpha=\frac{-y\,\dd x+x\,\dd y}{x^2+y^2}.
\]
它正是前一节极坐标中写出的局部 \(\dd\theta\)，
但右侧本身在整个穿孔平面有定义。
逐项求导得 \(\dd\alpha=0\)，所以它处处局部有势。
但沿单位圆 \((x,y)=(\cos t,\sin t)\) 有 \(\alpha=\dd t\)，
故 \(\oint\alpha=2\pi\)。若 \(\alpha=\dd f\) 且 \(f\) 是全局
单值函数，链式法则使闭路积分等于 \(f(\gamma(2\pi))-f(\gamma(0))=0\)，
矛盾。局部可积与全局单值是两个不同问题；
电磁势的规范和几何上的拓扑障碍都沿着这个区别发展。

\section{向量场的流与 李括号}
\label{sec:math-lie-bracket-flow}

切向量只属于一个点；向量场 \(X\) 给每个点分配一个切向量，
于是可沿其方向解常微分方程
\(\dot\gamma(t)=X(\gamma(t))\)。局部解的映射
\(\Phi_t(p)=\gamma_p(t)\) 称为 \(X\) 的流。
先沿 \(X\) 走一小段再沿 \(Y\) 走，与交换顺序一般不同。

\begin{definition}[向量场的李括号]
向量场 \(X,Y\) 都可作为函数上的微分算子。定义
\[
[X,Y](f)=X(Y(f))-Y(X(f)).
\]
这是另一个向量场，称李括号。
\end{definition}

为了核对它确为向量场，把 \(fg\) 代入上式并两次使用 Leibniz 法则；
含 \(X(f)Y(g)\) 与 \(Y(f)X(g)\) 的交叉项互相抵消，
剩下 \([X,Y](f)g+f[X,Y](g)\)。
由前一节“导子就是切向量”的定理，所得运算在每一点确是切向量。
若 \(X=X^i\partial_i,Y=Y^i\partial_i\)，对坐标函数作用得到
\[
[X,Y]=(X^j\partial_jY^i-Y^j\partial_jX^i)\partial_i.
\]
因此坐标向量场 \(\partial_i\) 两两对易。

\begin{definition}[标架与余标架]
开集上的局部标架是每一点的一组光滑切向量基
\(e_1,\ldots,e_n\)。满足 \(\theta^a(e_b)=\delta^a{}_b\)
的光滑一形式 \(\theta^1,\ldots,\theta^n\) 称为对偶余标架。
\end{definition}

标架可以不是坐标基。写
\([e_a,e_b]=C^c{}_{ab}e_c\)，
将外微分的一形式公式用于 \(e_a,e_b\)，得到
\[
\dd\theta^c(e_a,e_b)
=e_a[\theta^c(e_b)]-e_b[\theta^c(e_a)]
-\theta^c([e_a,e_b])
=-C^c{}_{ab}.
\]
因 \(C^c{}_{ab}=-C^c{}_{ba}\)，等价地
\[
\dd\theta^c=-\frac12C^c{}_{ab}\theta^a\wedge\theta^b.
\]
这条式子只测量标架是否来自坐标，尚未测量空间曲率。

在平面极坐标区域 \(r>0\) 取正交标架
\(e_r=\partial_r,\ e_\theta=r^{-1}\partial_\theta\)。
直接用上面的坐标表达算出
\([e_r,e_\theta]=-r^{-1}e_\theta\)。
其余标架是 \(\theta^r=\dd r,\theta^\theta=r\,\dd\theta\)，
所以 \(\dd\theta^\theta=\dd r\wedge \dd\theta
=r^{-1}\theta^r\wedge\theta^\theta\)。
平面仍是平的；非零括号源于这组正交基随位置转动。




\section{流与外微分的公式联系}
流、括号和外微分在张量的输运中同时出现。下面把已经证明的恒等式并列起来，便于辨认它们的作用对象。
\begin{note}[公式之间的联系]
\begin{align*}
\dd^2&=0,\\
F^*\dd&=\dd F^*,\\
[X,Y](f)&=X(Yf)-Y(Xf),\\
\mathcal L_Xf&=X(f),\\
\mathcal L_XY&=[X,Y],\\
\mathcal L_X\alpha&=\iota_{X}\dd\alpha+\dd\iota_{X}\alpha.
\end{align*}
主要定义与公式见\eqnrefs{eq:math-dg-differential-function}、\eqnrefs{eq:cartan-lie-lie}、\eqnrefs{eq:math-dg-synthesis-lie-flow}与\eqnrefs{eq:math-dg-synthesis-cartan}。
\end{note}

\begin{proposition}[切向量的两种定义]
设 $M$ 为光滑流形、$p\in M$。通过 $p$ 的光滑曲线按相同初速度取商所得空间，与点 $p$ 处满足 Leibniz 规则的实线性导子空间自然同构。对应为 $[\gamma]\mapsto[f\mapsto (f\circ\gamma)^\prime(0)]$。
\end{proposition}
\begin{proof}
需证两个方向：每个曲线速度是导子，每个导子由某曲线速度实现。

方向一：曲线速度 $\Rightarrow$ 导子。任给光滑曲线 $\gamma:(-\varepsilon,\varepsilon)\to M$，$\gamma(0)=p$，定义 $v_\gamma(f)=\left.\frac{\dd}{\dd t}\right|_0 f(\gamma(t))$。对 $f,g\in C^\infty(M)$，由链式法则
\begin{align*}
 v_\gamma(fg)
 &=\left.\frac{\dd}{\dd t}\right|_0\bigl(f(\gamma(t))\,g(\gamma(t))\bigr)\\
 &=\dot f(\gamma(0))\,g(\gamma(0))+f(\gamma(0))\,\dot g(\gamma(0))\\
 &=v_\gamma(f)\,g(p)+f(p)\,v_\gamma(g),
\end{align*}
故 $v_\gamma$ 满足 Leibniz 规则，是 $T_pM$ 中的导子。线性性 $v_\gamma(af+bg)=a\,v_\gamma(f)+b\,v_\gamma(g)$ 直接由导数线性得到。

方向二：导子 $\Rightarrow$ 曲线速度。先说明导子的坐标表示没有额外假设。Leibniz 规则给 $v(1)=0$。若函数 $f$ 在 $p$ 附近恒零，取邻域内等于一且支撑在该邻域的光滑截断函数 $\chi$，则 $\chi f=0$ 且
\[
v(\chi f)=v(\chi)f(p)+\chi(p)v(f)=v(f)=0.
\]
因此 $v$ 只依赖函数在 $p$ 的芽。坐标邻域内的 Hadamard 分解给
\[
f(x)-f(p)=\sum_i(x^i-x^i(p))h_i(x),\qquad h_i(p)=\partial_i f(p).
\]
对它作用 $v$，由 Leibniz 规则得到 $v(f)=\sum_i v(x^i)\partial_i f(p)$。

给定 $v\in T_pM$（导子定义），在局部坐标 $(x^1,\ldots,x^n)$ 中写
\begin{equation*}
 v=v^i\,\partial_i\big|_p,
 \qquad
 v^i=v(x^i).
\end{equation*}
构造坐标曲线 $\gamma:(-\varepsilon,\varepsilon)\to M$，令
\begin{equation*}
 x^i(\gamma(t))=x^i(p)+t\,v^i,
 \qquad i=1,\ldots,n,
\end{equation*}
在足够小 $\varepsilon$ 上限制于坐标邻域内。对任意 $f\in C^\infty(M)$，在坐标中写 $f=f(x^1,\ldots,x^n)$，则
\begin{align*}
 \left.\frac{\dd}{\dd t}\right|_0 f(\gamma(t))
 &=\left.\frac{\dd}{\dd t}\right|_0 f\bigl(x^1(p)+tv^1,\ldots,x^n(p)+tv^n\bigr)\\
 &=\sum_{i=1}^n \frac{\partial f}{\partial x^i}\bigg|_p v^i
 =v^i\,\partial_i f(p)=v(f).
\end{align*}
因此 $v_\gamma=v$，每个导子都由某条曲线速度实现。

若两条曲线 $\gamma_1,\gamma_2$ 给出同一导子，则对所有 $f$ 有 $v_{\gamma_1}(f)=v_{\gamma_2}(f)$，特别取 $f=x^i$ 得 $\dot x^i_1(0)=\dot x^i_2(0)$，即速度分量相同。反之，速度分量相同则对任意 $f$ 的一阶作用相同。因此映射 $[\gamma]\mapsto v_\gamma$ 是曲线等价类到导子的双射，且与坐标选取无关，两种定义自然同构。
\end{proof}

\begin{lemma}[Poincaré 引理与缩放同伦]
设 $U\subset\RR^n$ 关于原点星形。对 $k\ge1$ 定义 $(K\alpha)_x(v_1,\ldots,v_{k-1})=\int_0^1t^{k-1}\alpha_{tx}(x,v_1,\ldots,v_{k-1})\,\dd t$。则 $\dd K+K\dd=\id-F_0^*$；特别地，每个正次数闭形式 $\alpha$ 都满足 $\alpha=\dd K\alpha$。
\end{lemma}
\begin{proof}
用缩放同伦 $F_t(x)=tx$ 将 $t=1$ 的形式 $\alpha$ 连续拉回到 $t=0$（原点），沿途的"变化率"由 Cartan 公式控制，积分后得到同伦恒等式。

对 $k$-形式 $\alpha$，定义
\begin{equation*}
 H:[0,1]\times U\to U,
 \qquad H(t,x)=tx.
\end{equation*}
沿 $t$ 方向的生成向量场在 $U$ 中对应径向场 $\mathcal R/t$（因 $\frac{\dd}{\dd t}(tx)=x=\mathcal R/t$ 在 $tx$ 处取值）。

由 Lie 导数与外微分的关系（Cartan 公式）$\mathcal L_X=\dd\iota_X+\iota_X\dd$ 以及拉回与 Lie 导数的交换关系，
\begin{align*}
 \frac{\dd}{\dd t}F_t^*\alpha
 &=F_t^*(\mathcal L_{\mathcal R/t}\alpha)
 =F_t^*\left(\dd\iota_{\mathcal R/t}\alpha+\iota_{\mathcal R/t}\dd\alpha\right)\\
 &=\frac{1}{t}F_t^*\left(\dd\iota_{\mathcal R}\alpha+\iota_{\mathcal R}\dd\alpha\right).
\end{align*}

从 $t=0$ 积分到 $1$：
\begin{equation*}
 \alpha-F_0^*\alpha
 =\int_0^1\frac{\dd}{\dd t}F_t^*\alpha\,\dd t
 =\dd\!\int_0^1 F_t^*(\iota_{\mathcal R}\alpha)\,\frac{\dd t}{t}
 +\int_0^1 F_t^*(\iota_{\mathcal R}\dd\alpha)\,\frac{\dd t}{t}.
\end{equation*}
拉回包含切向量的缩放，不能把径向向量在 $tx$ 处的值再重复乘一次 $t$。对 $k$-形式取分量后得到 $t^{k-1}\alpha_{tx}(x,\cdot)$；不要求 $\alpha$ 本身齐次。记
\begin{equation*}
 (K\alpha)_x=\int_0^1 F_t^*(\iota_{\mathcal R}\alpha)\,\frac{\dd t}{t},
\end{equation*}
即\eqnrefs{eq:poincare-homotopy-operator}，则上式化为
\begin{equation*}
 \alpha-F_0^*\alpha=\dd K\alpha+K\dd\alpha,
 \end{equation*}
也就是\eqnrefs{eq:poincare-homotopy-identity}。

若 $k\ge1$ 且 $\dd\alpha=0$（闭形式），常值映射 $F_0$ 把 $k$-形式拉回为零（$k\ge1$ 时原点处切空间退化），故 $F_0^*\alpha=0$。代入同伦恒等式：
\begin{equation*}
 \alpha=\dd(K\alpha),
\end{equation*}
即 $\alpha$ 恰当，势为 $K\alpha$。

这说明局部上"闭则恰当"来自空间可沿径向连续缩到一点；全局上同调正是这种收缩在整体上失败留下的信息。
\end{proof}

\begin{proposition}[正则水平集的切空间]
若 $F:M\to N$ 光滑、$y\in N$ 是正则值，则 $F^{-1}(y)$ 是余维 $\dim N$ 的嵌入子流形，且对每个 $p\in F^{-1}(y)$，$T_pF^{-1}(y)=\ker F_{*p}$。
\end{proposition}
\begin{proof}
分两步证明双向包含。先证 $T_pF^{-1}(y)\subseteq\ker F_{*p}$。任取切向量 $v\in T_pF^{-1}(y)$，由切空间定义存在光滑曲线 $\gamma:(-\delta,\delta)\to F^{-1}(y)$ 满足 $\gamma(0)=p$、$\dot\gamma(0)=v$。因 $\gamma$ 像落在水平集中，
\[
F(\gamma(t))\equiv y
\]
对一切 $t$ 成立。对 $t$ 在 $t=0$ 处求导，由链式法则
\[
F_{*p}\dot\gamma(0)
=\frac{\dd}{\dd t}\bigg|_{t=0}F(\gamma(t))
=\frac{\dd}{\dd t}\bigg|_{t=0}y=0,
\]
即 $F_{*p}v=0$，故 $v\in\ker F_{*p}$。

再证 $\ker F_{*p}\subseteq T_pF^{-1}(y)$。因 $y$ 是正则值，$F_{*p}$ 满射，由淹没正规形可取 $p$ 附近局部坐标使
\[
F(x^1,\ldots,x^m)=(x^1,\ldots,x^n).
\]
在此坐标下，水平集 $F^{-1}(y)$ 局部由
\[
x^1=y^1,\;\ldots,\;x^n=y^n
\]
定义，其余坐标 $x^{n+1},\ldots,x^m$ 自由。因此 $F^{-1}(y)$ 可参数化为
\[
(t^{n+1},\ldots,t^m)\longmapsto
(y^1,\ldots,y^n,t^{n+1},\ldots,t^m),
\]
其切空间由 $\partial_{n+1},\ldots,\partial_m$ 张成。而
\[
F_{*p}(\partial_a)=\partial_a\quad(a\le n),
\qquad
F_{*p}(\partial_a)=0\quad(a>n),
\]
故 $\ker F_{*p}=\operatorname{span}\{\partial_{n+1},\ldots,\partial_m\}=T_pF^{-1}(y)$。综合双向包含得到\eqnrefs{eq:tangent-regular-level-set}。
\end{proof}

\begin{lemma}[拉回与外微分交换]
对任意光滑映射 $F:M\to N$ 与微分形式 $\alpha\in\Omega^k(N)$，有 $\dd(F^*\alpha)=F^*(\dd\alpha)$。
\end{lemma}
\begin{proof}
目标是证明拉回与外微分可交换，即 $\dd F^*=F^*\dd$。在 $N$ 上取局部坐标 $y^a$，把任意 $k$-形式写成
\[
\alpha
=\frac1{k!}\alpha_{a_1\cdots a_k}(y)
\dd y^{a_1}\wedge\cdots\wedge\dd y^{a_k}.
\]
记 $F^a:=y^a\circ F$。由拉回定义，
\[
F^*\alpha
=\frac1{k!}(\alpha_{a_1\cdots a_k}\circ F)
\dd F^{a_1}\wedge\cdots\wedge\dd F^{a_k}.
\]
对它作用 $\dd$。由于 $\dd^2F^a=0$，外微分只作用在系数上：
\[
\begin{aligned}
\dd(F^*\alpha)
={}&\frac1{k!}
\dd(\alpha_{a_1\cdots a_k}\circ F)
\wedge\dd F^{a_1}\wedge\cdots\wedge\dd F^{a_k}\\
={}&\frac1{k!}
\bigl[(\partial_b\alpha_{a_1\cdots a_k})\circ F\bigr]
\dd F^b\wedge\dd F^{a_1}\wedge\cdots\wedge\dd F^{a_k}.
\end{aligned}
\]
另一方面，先在 $N$ 上计算
\[
\dd\alpha
=\frac1{k!}\partial_b\alpha_{a_1\cdots a_k}
\dd y^b\wedge\dd y^{a_1}\wedge\cdots\wedge\dd y^{a_k},
\]
再拉回得到完全相同的式子。因此 $\dd F^*=F^*\dd$。证明只使用链式法则与 $\dd^2=0$，所以该恒等式与任何度量或联络无关。

取 $k=1$ 且 $F(t)=(\cos t,\sin t)$ 可显式验证两端相等。
\end{proof}

\begin{theorem}[横截原像定理]
设 $S\subset N$ 是嵌入子流形，光滑映射 $F:M\to N$ 横截于 $S$。则 $F^{-1}(S)$ 为余维 $\operatorname{codim}_N S$ 的嵌入子流形，且 $T_pF^{-1}(S)=\{v\in T_pM:F_{*p}v\in T_{F(p)}S\}$。
\end{theorem}
\begin{proof}
在 $q=F(p)\in S$ 附近选适配坐标 $(y^1,\ldots,y^r,y^{r+1},\ldots,y^n)$，使
\begin{equation*}
 S\cap U=\{y^{r+1}=\cdots=y^n=0\},
 \qquad r=\dim S.
\end{equation*}
即 $S$ 的前 $r$ 个坐标自由，后 $n-r$ 个坐标为零。定义法向投影
\begin{equation*}
 \pi_\perp:N\supset U\to\RR^{n-r},
 \qquad
 \pi_\perp(y)=(y^{r+1},\ldots,y^n),
\end{equation*}
它提取后 $n-r$ 个"法向"坐标分量。

在此坐标下，
\begin{equation*}
 F^{-1}(S)=\{p\in M:F(p)\in S\}
 =\{p\in M:\pi_\perp(F(p))=0\}
 =(\pi_\perp\circ F)^{-1}(0).
\end{equation*}
因此 $F^{-1}(S)$ 是复合映射 $\pi_\perp\circ F:M\to\RR^{n-r}$ 的零水平集。

$\pi_\perp$ 的微分是到后 $n-r$ 个坐标分量的投影。复合映射的微分
\begin{equation*}
 \dd(\pi_\perp\circ F)_p=\pi_\perp\circ\dd F_p:T_pM\to\RR^{n-r}
\end{equation*}
满射，当且仅当 $\dd F_p$ 的像加上 $S$ 的切空间（即前 $r$ 个坐标方向）张成整个 $T_{F(p)}N$，这正是横截条件\eqnrefs{eq:math-dg-transversality}。

既然 $0$ 是 $\pi_\perp\circ F$ 的正则值，正则值定理给出 $F^{-1}(S)=(\pi_\perp\circ F)^{-1}(0)$ 是 $M$ 的嵌入子流形，余维数 $n-r=\operatorname{codim}_N S$，即式\eqnrefs{eq:math-dg-transverse-codimension}。切空间是 $\dd(\pi_\perp\circ F)_p$ 的核，即 $\{v\in T_pM:\dd F_p(v)\in T_{F(p)}S\}$，即\eqnrefs{eq:math-dg-transverse-tangent-space}。

横截原像定理因而不是新的局部正规形，而是把正则值定理应用于法向坐标分量。
\end{proof}

\begin{proposition}[相关向量场的流]
若 $F:M\to N$ 光滑且 $F_*X=Y\circ F$，则在两边均有定义的局部流区间内，$F\circ\Phi_t^X=\Phi_t^Y\circ F$。当 $F$ 为微分同胚时，这就是两流的共轭关系。
\end{proposition}
\begin{proof}
分两步证明 $F\circ\Phi_t^X=\Phi_t^Y\circ F$。

固定 $p\in M$，令
\[
 \eta(t)=F(\Phi_t^X(p)).
\]
这是 $F$ 沿 $X$ 的流从 $p$ 出发的轨道所得到的 $N$ 上的曲线，满足初始条件 $\eta(0)=F(p)$。

对 $\eta(t)=F(\Phi_t^X(p))$ 关于 $t$ 求导。由链式法则，
\[
 \dot\eta(t)
 =F_{*\Phi_t^X(p)}\!\left(\frac{\dd}{\dd t}\Phi_t^X(p)\right)
 =F_{*\Phi_t^X(p)}\,X_{\Phi_t^X(p)}.
\]
由相关向量场条件 $F_*X=Y\circ F$，在点 $\Phi_t^X(p)$ 处，
\[
 F_{*\Phi_t^X(p)}\,X_{\Phi_t^X(p)}
 =Y_{F(\Phi_t^X(p))}
 =Y_{\eta(t)}.
\]
因此 $\dot\eta(t)=Y_{\eta(t)}$，即 $\eta$ 是 $Y$ 的积分曲线。

另一方面，$\Phi_t^Y(F(p))$ 也是 $Y$ 的积分曲线，且满足同一初始条件：
\[
 \eta(0)=F(p)=\Phi_0^Y(F(p)).
\]
由 ODE 解的唯一性（在流的存在区间内），
\[
 \eta(t)=\Phi_t^Y(F(p)),
\]
即 $F(\Phi_t^X(p))=\Phi_t^Y(F(p))$。因 $p$ 任意，得到\eqnrefs{eq:related-flow-naturality}。
\end{proof}

\begin{proposition}[流交换子的混合二阶项]
设 $X,Y$ 为光滑向量场，流分别为 $\Phi_t,\Psi_s$。定义 $C_{t,s}=\Psi_{-s}\circ\Phi_{-t}\circ\Psi_s\circ\Phi_t$，则 $C_{t,s}^*f=f+ts[X,Y]f+O(|t|^2|s|+|t||s|^2)$。相反次序 $\Phi_t\circ\Psi_s\circ\Phi_{-t}\circ\Psi_{-s}$ 的混合项为 $-ts[X,Y]$。
\end{proposition}
\begin{proof}
将 $C_{t,s}=\Psi_{-s}\circ\Phi_{-t}\circ\Psi_s\circ\Phi_t$ 的拉回展开到 $ts$ 阶，证明最低非平凡项恰好是 Lie 括号 $[X,Y]$。

令 $\Phi_t$ 与 $\Psi_s$ 分别是 $X$ 与 $Y$ 的局部流。对任意函数 $f$，由流定义，
\begin{equation*}
 \Phi_t^*f=f+t\,Xf+O(t^2),
 \qquad
 \Psi_s^*f=f+s\,Yf+O(s^2).
\end{equation*}
等价地，拉回算子的一阶展开为
\begin{equation*}
 \Phi_{\pm t}^*=I\pm t\,X+O(t^2),
 \qquad
 \Psi_{\pm s}^*=I\pm s\,Y+O(s^2).
\end{equation*}

考虑群交换子对应的拉回组合
\begin{equation*}
 \mathcal C_{t,s}^*
 :=\Phi_t^*\Psi_s^*\Phi_{-t}^*\Psi_{-s}^*.
\end{equation*}
依次代入一阶展开，并只保留 $ts$ 阶交叉项（纯 $t$ 项与纯 $s$ 项由逆流精确抵消，余项至少含一个 $t$ 和一个 $s$）：
\begin{align*}
 \mathcal C_{t,s}^*
 &=(I+tX)(I+sY)(I-tX)(I-sY)+O(t^2s,ts^2)\\
 &=I+ts(XY-YX)+O(t^2s,ts^2)\\
 &=I+ts[X,Y]+O(t^2s,ts^2).
\end{align*}
混合项的逐项核对：前两个因子给 $tsXY$，后两个因子也给 $tsXY$；两个乘积之间的交叉项给 $-tsXY-tsYX$。合计为 $ts(XY-YX)$。纯二阶项必须连同每个流的二阶 Taylor 项一起处理，不能仅凭四个一阶多项式就声称它们消失。

这里拉回逆转映射复合的次序。若直接采用 $\Phi_t\circ\Psi_s\circ\Phi_{-t}\circ\Psi_{-s}$，其拉回依次为 $\Psi_{-s}^*\Phi_{-t}^*\Psi_s^*\Phi_t^*$，混合项因而为 $-ts[X,Y]$。

Lie 括号不是人为附加的代数运算，而是两个局部流不交换性的最低非平凡系数：若 $[X,Y]=0$，群交换子在 $ts$ 阶消失，流可交换（Frobenius 可积条件的无穷小版本）。符号由组合次序完全固定；反转群交换子次序，括号也相应变号。
\end{proof}

\begin{proposition}[Cartan 演算的闭合关系]
对光滑向量场 $X,Y$，在整个微分形式代数上有 $[\mathcal L_X,\iota_Y]=\iota_{[X,Y]}$、$[\mathcal L_X,\mathcal L_Y]=\mathcal L_{[X,Y]}$ 和 $[\mathcal L_X,\dd]=0$。插入算子满足 $[\iota_X,\iota_Y]_{\rm gr}=0$。
\end{proposition}
\begin{proof}
先在函数与一形式上验证
\[
[\mathcal L_X,\iota_Y]=\iota_{[X,Y]}.
\]
对一形式 $\alpha$，
\[
\begin{aligned}
([\mathcal L_X,\iota_Y]\alpha)
&=X(\alpha(Y))-(\mathcal L_X\alpha)(Y)\\
&=X(\alpha(Y))-
\bigl[X(\alpha(Y))-\alpha([X,Y])\bigr]\\
&=\alpha([X,Y]).
\end{aligned}
\]
两边都是次数 $-1$ 的导子，所以在整个外代数上一致，得到式\eqnrefs{eq:cartan-lie-iota}。再用分次 Jacobi 恒等式与
\[
\mathcal L_X=[\dd,\iota_X]_{\rm gr}
\]
可得
\[
\begin{aligned}
[\mathcal L_X,\mathcal L_Y]
&=[\mathcal L_X,[\dd,\iota_Y]_{\rm gr}]\\
&=[\dd,[\mathcal L_X,\iota_Y]]_{\rm gr}\\
&=[\dd,\iota_{[X,Y]}]_{\rm gr}\\
&=\mathcal L_{[X,Y]},
\end{aligned}
\]
即\eqnrefs{eq:cartan-lie-lie}。所以 Lie 括号、流与微分形式代数三者并不是并列工具，而是由同一分次交换子结构闭合起来。

在 $\RR^3$ 上显式验证 Cartan 恒等式。取 $X=\partial_x$，$\alpha=\dd y$。则 $\iota_X\alpha=0$，$\mathcal L_X\alpha=\dd(\iota_X\alpha)+\iota_X\dd\alpha=0$（因 $\dd\alpha=0$）。另一方向，$\mathcal L_X\dd y=\dd(\mathcal L_X y)=\dd(0)=0$，验证 $[\mathcal L_X,\dd]=0$。取 $Y=x\partial_y$，$[X,Y]=\partial_y$，则 $[\mathcal L_X,\iota_Y]\dd y=\iota_{[X,Y]}\dd y=1$，直接计算 $\mathcal L_X(\dd y(Y))-\iota_Y(\mathcal L_X\dd y)=\mathcal L_X(x)-0=1$，验证\eqnrefs{eq:cartan-lie-iota}。

这一分次微分结构在多处反复出现：Bianchi 恒等式 $\dd F=\dd^2A=0$ 给出齐次 Maxwell 方程，不可压 Euler 的涡度方程 $\partial_t\omega+\mathcal L_v\omega=0$ 是 Cartan 恒等式的特例，Maurer--Cartan 方程则是 $\dd^2=0$ 的非线性推广。

故外微分与 Hodge 对偶把电磁场强、Bianchi 恒等式与 Maxwell 方程 $\dd *F=J$ 写成统一语言，其核心 $\dd^2=0$ 与 Cartan 恒等式根源即外代数上的分次微分结构。
\end{proof}

\begin{theorem}[Cartan 恒等式]
对任意光滑向量场 $X$ 与任意次数的微分形式，Lie 导数满足 $\mathcal L_X=\iota_X\dd+\dd\iota_X$。该恒等式不需要度量或联络。
\end{theorem}
\begin{proof}
目标是证明 Cartan 恒等式 $\mathcal L_X=\iota_X\dd+\dd\iota_X$。定义次数为零的算子
\begin{equation*}
\mathscr D_X:=\mathcal L_X-\iota_{X}\dd-\dd\iota_{X}.
\end{equation*}
对任意函数 $f$，由 $\iota_{X}f=0$ 与\eqnrefs{eq:math-dg-differential-function}可得
\begin{align*}
\mathscr D_Xf
&=\mathcal L_X f-\iota_{X}\dd f-\dd(\iota_{X}f)\\
&=X(f)-\iota_{X}\dd f-\dd 0\\
&=X(f)-\dd f(X)=0.
\end{align*}
Lie 导数是次数零导子，而 $\iota_{X}\dd+\dd\iota_{X}$ 也是次数零导子：对任意形式 $\omega,\eta$，
\begin{align*}
(\iota_{X}\dd+\dd\iota_{X})(\omega\wedge\eta)
&=\iota_{X}\dd(\omega\wedge\eta)+\dd\iota_{X}(\omega\wedge\eta)\\
&=\iota_{X}(\dd\omega\wedge\eta+(-1)^{|\omega|}\omega\wedge\dd\eta)\\
&\quad+\dd(\iota_{X}\omega\wedge\eta+(-1)^{|\omega|}\omega\wedge\iota_{X}\eta),
\end{align*}
按 $\iota_{X}$ 与 $\dd$ 各自的反导子 Leibniz 规则展开后，非齐次项相消，剩余项恰为
\begin{equation*}
(\iota_{X}\dd+\dd\iota_{X})\omega\wedge\eta
+\omega\wedge(\iota_{X}\dd+\dd\iota_{X})\eta.
\end{equation*}
因此 $\mathscr D_X$ 满足普通 Leibniz 规则。又由\eqnrefs{eq:math-dg-synthesis-lie-d-commute}和 $\dd^2=0$，
\begin{align*}
[\mathscr D_X,\dd]
&=[\mathcal L_X,\dd]-[\iota_{X}\dd,\dd]-[\dd\iota_{X},\dd]\\
&=0-\iota_{X}\dd^2+\dd\iota_{X}\dd-\dd\iota_{X}\dd+\dd^2\iota_{X}=0.
\end{align*}
于是对任意函数 $f$，
\begin{equation*}
\mathscr D_X(\dd f)=\dd(\mathscr D_X f)=\dd 0=0.
\end{equation*}
局部上任意微分形式都由光滑函数与坐标一形式 $\dd x^i$ 通过外积生成，而 $\mathscr D_X$ 既是导子又与 $\dd$ 交换，故对单项式 $f_0\,\dd f_1\wedge\cdots\wedge\dd f_k$ 有
\begin{align*}
\mathscr D_X(f_0\,\dd f_1\wedge\cdots\wedge\dd f_k)
&=(\mathscr D_X f_0)\,\dd f_1\wedge\cdots\wedge\dd f_k\\
&\quad+\sum_{j=1}^k f_0\,\dd f_1\wedge\cdots\wedge\mathscr D_X(\dd f_j)\wedge\cdots\wedge\dd f_k\\
&=0.
\end{align*}
所以导子 $\mathscr D_X$ 在整个 $\Omega^\bullet(M)$ 上恒为零。故
\begin{equation}\label{eq:math-dg-synthesis-cartan}
{\mathcal L_X=\iota_{X}\dd+\dd\iota_{X}.}
\end{equation}
\end{proof}

